% संपूर्ण मराठी ग्रंथातील प्रकरण 56: प्रस्तावना
% हा संपादनयोग्य स्रोतखंड आहे. संकलनासाठी संपूर्ण स्रोतसंग्रहातील
% अक्षररूपे, पूर्वभाग, चिन्हव्याख्या आणि आकृती-स्रोत आवश्यक आहेत.
% मूळ: Open Logic Project; मजकूर आणि रूपांतर CC BY 4.0.
% यंत्रानुवाद, दुरुस्ती आणि तपासणी: OpenAI Codex —
% GPT-5.6 Sol आणि GPT-6 Sol, दोन्ही Ultra effort.
% दुरुस्ती-पुनरावलोकन: GPT-6.1 Sol, Ultra effort.
\chapter{प्रस्तावना}\label{cnt:int:chap}









\section{वास्तविक अभिव्यंजन}\label{cnt:int:mat:sec}

इंग्रजीत सशर्त विधानाचे सर्वात सोपे रूप
‘‘जर \dots तर \dots’’ असे असते; येथे
\dots{} यांच्या जागी स्वतः विधाने येतात,
उदाहरणार्थ, ‘‘जर हे काम सेवकाने केले
असेल, तर माळी निर्दोष आहे.’’ प्राथमिक
तर्कशास्त्रात अशी सशर्त विधाने
$\lif$ या कारकाने दर्शवायला शिकतो:
\dots{} ने दाखवलेले भाग, उदाहरणार्थ,
सूत्रे $\varphi$ आणि~$\psi$ यांनी
दर्शवायचे; मग संपूर्ण सशर्त विधान
$\varphi \lif \psi$ असे दर्शवतात.

$\lif$ हा कारक \emph{सत्यता-फलनात्मक}
आहे: $\varphi \lif \psi$ चे सत्यतामूल्य---$\True$
किंवा $\False$---हे $\varphi$ आणि~$\psi$
यांच्या सत्यतामूल्यांवरून ठरते.
$\varphi$ असत्य असेल किंवा $\psi$ सत्य
असेल तेव्हा आणि केवळ तेव्हाच $\varphi \lif \psi$
सत्य असते; अन्यथा असत्य असते.
$\pAssign v$ या सत्यतामूल्य-नियुक्तीच्या
संदर्भात $\pSat{v}{\varphi \lif \psi}$
तेव्हा आणि केवळ तेव्हाच, जेव्हा
$\pSat/{v}{\varphi}$ किंवा
$\pSat{v}{\psi}$, अशी व्याख्या
करतो. ही चिन्हार्थमीमांसा
असलेल्या $\lif$ कारकाला
\emph{वास्तविक अभिव्यंजन}
म्हणतात.

या व्याख्येवरून काही प्राथमिक तार्किक
तथ्ये मिळतात. सर्वप्रथम,
वास्तविक अभिव्यंजनासाठी निगमन प्रमेय लागू होते:
\begin{align}
  \text{जर } \Gamma, \varphi \Entails \psi \text{ तर } \Gamma \Entails \varphi \lif \psi
\end{align}
ते सत्यता-फलनात्मक आहे:
$\varphi \lif \psi$ आणि
$\lnot \varphi \lor \psi$ सममूल्य आहेत:
\begin{align}
  \varphi \lif \psi & \Entails \lnot \varphi \lor \psi\\
  \lnot \varphi \lor \psi & \Entails \varphi \lif \psi
  \intertext{वास्तविक अभिव्यंजन त्याच्या
    उत्तरांगावरून आणि पूर्वांगाच्या
    निषेधावरून तार्किकरीत्या निष्पन्न होते:}
  \psi & \Entails \varphi \lif \psi\\
  \lnot \varphi & \Entails \varphi \lif \psi
  \intertext{असत्य वास्तविक अभिव्यंजन
    हे त्याच्या पूर्वांगाचा आणि उत्तरांगाच्या
    निषेधाचा संधियोग याच्याशी सममूल्य असते:
    $\varphi \lif \psi$ असत्य असेल, तर
    $\varphi \land \lnot \psi$ सत्य असते,
    आणि उलटही:}
  \lnot(\varphi \lif \psi) & \Entails \varphi \land \lnot \psi\\
  \varphi \land \lnot \psi & \Entails \lnot(\varphi \lif \psi)
  \intertext{वास्तविक अभिव्यंजनामुळे
    विध्यात्मक अनुमान लागू होते:}
  \varphi, \varphi \lif \psi & \Entails \psi
  \intertext{वास्तविक अभिव्यंजन
    निष्कर्षांचे संयुग्मीकरण करते:}
  \varphi \lif \psi,  \varphi \lif \chi & \Entails \varphi \lif (\psi \land \chi)
  \intertext{पूर्वांग अधिक बलवान
    नेहमी करता येते; म्हणजे
    हे अभिव्यंजन \emph{एकस्वनिक} आहे:}
  \varphi \lif \psi & \Entails (\varphi \land \chi) \lif \psi
  \intertext{वास्तविक अभिव्यंजन
    संक्रमक आहे; म्हणजे साखळी-नियम वैध आहे:}
  \varphi \lif \psi, \psi \lif \chi & \Entails \varphi \lif \chi
  \intertext{वास्तविक अभिव्यंजन
    त्याच्या प्रतिपरिवर्तित विधानाशी सममूल्य आहे:}
  \varphi \lif \psi & \Entails \lnot \psi \lif \lnot \varphi\\
  \lnot \psi \lif \lnot \varphi & \Entails \varphi \lif \psi
\end{align}

हे सर्व निष्कर्ष गणिती युक्तिवादात
उपयुक्त आणि निर्विवाद आहेत.
मात्र गणिताबाहेरील संदर्भांतील
सममूल्य निष्कर्षांना कथित
प्रतिदृष्टांत तत्त्वज्ञान आणि भाषाविज्ञान
यांच्या साहित्यात विपुल आहेत.
त्यांवरून असे सुचते की
इंग्रजीतील ‘‘जर \dots तर \dots’’
विधाने दर्शवताना
वास्तविक अभिव्यंजन~$\lif$
हा कारक योग्य नसतो---किंवा
किमान नेहमीच योग्य नसतो.












\section{वास्तविक अभिव्यंजनाचे विरोधाभास}\label{cnt:int:par:sec}

इंग्रजीतील ‘‘जर \dots तर \dots’’
विधानांचे प्रतीकीकरण $\varphi \lif \psi$
या रूपात करण्यावर टीका करणाऱ्यांपैकी
C.~I. Lewis हे सुरुवातीचे एक होते.
व्हाइटहेड आणि रसेल यांच्या
\emph{Principia Mathematica} मध्ये
वास्तविक अभिव्यंजनाचा वापर करताना
$\lif$ चा उच्चार ‘‘यावरून
तार्किकरीत्या निष्पन्न होते’’ असा
केला जात असे; त्यावर लुईस यांची
टीका होती. $\lif$ चा अर्थ
‘‘तार्किकरीत्या निष्पन्न होते’’
असा घेतला, तर कोणतेही असत्य
विधान~$p$ यावरून $p$ यावरून~$q$
निष्पन्न होते असे ठरेल, कारण
$p$ असत्य असताना
$p \lif (p \lif q)$ सत्य असते.
त्याचप्रमाणे कोणतेही सत्य
विधान~$q$ यावरून $p$ यावरून~$q$
निष्पन्न होते असे ठरेल, कारण
$q$ सत्य असताना
$q \lif (p \lif q)$ सत्य असते,
ही लुईस यांची रास्त तक्रार होती.

तर्कशास्त्रज्ञांना अर्थातच माहीत
आहे की \emph{तार्किक निष्पन्नता}
हा कारक नाही; तो सूत्रे
किंवा विधानांमधील संबंध आहे.
म्हणून गोंधळ टाळण्यासाठी
$\lif$ ला ‘‘तार्किकरीत्या निष्पन्न
होते’’ असे वाचूच नये.\footnote{
  गणितज्ञ आणि संगणकशास्त्रज्ञ
  ‘‘$\lif$’’ चा उच्चार ‘‘यावरून
  निष्पन्न होते’’ असा अजूनही
  वारंवार करतात; पण लुईस यांनी
  दाखवलेला गोंधळ टाळण्यासाठी
  तत्त्वज्ञ ‘‘तरच’’ असा
  उच्चार करतात.} तसे
वाचत नसू, तर लुईस यांची
ती विशिष्ट चिंता उरत नाही:
$p$ हे एखादे असत्य इंग्रजी
विधान दर्शवत असले, तरी
$p$ यावरून $q$ तार्किकरीत्या
‘‘निष्पन्न’’ होत नाही.
$p \Entails q$ आहे का हे
ठरवण्यासाठी \emph{सर्व}
सत्य-मूल्यांकने विचारात घ्याव्या
लागतात; योगायोगाने असत्य
ठरणारे विधान दर्शवण्यासाठी
$p$ वापरले, तरी
$p \Entails/ q$.

तरीही लुईस यांच्या पुढीलसारख्या
‘‘जर \dots तर\dots’’ विधानात
काहीतरी विचित्र वाटते:
\begin{quote}
चंद्र हिरव्या चीजचा बनलेला असेल,
तर $2+2=4$.
\end{quote}
आणि पुढील अनुमानांतही:
\begin{quote}
  चंद्र हिरव्या चीजचा बनलेला नाही.
  म्हणून चंद्र हिरव्या चीजचा
  बनलेला असेल, तर $2+2=4$.

  $2+2 = 4$. म्हणून चंद्र
  हिरव्या चीजचा बनलेला
  असेल, तर $2+2=4$.
\end{quote}
पण ‘‘जर \dots तर \dots’’
याचा अर्थ केवळ $\lif$
असा असेल, तर ते विधान
निर्विवाद सत्य आणि ही
अनुमाने निर्विवाद वैध ठरतील.

दुसरा प्रश्न $(\varphi \lif \psi) \lor (\psi \lif
\varphi)$ या सर्वतः सत्य सूत्राबद्दल
आहे. त्यावरून असे सुचेल की
वृत्तपत्रातून $S$ आणि $T$
ही कोणतीही दोन निश्चयार्थी
वाक्ये निवडली, तर ‘‘जर
$S$ तर $T$, किंवा जर
$T$ तर~$S$’’ हे विधान
सत्य असले पाहिजे.












\section{कठोर अभिव्यंजन}\label{cnt:int:str:sec}

लुईस यांनी \emph{कठोर अभिव्यंजन}~$\strictif$
हा कारक मांडला आणि वास्तविक अभिव्यंजनाऐवजी
तो तार्किक निष्पन्नतेशी जुळतो, असा युक्तिवाद
केला. अवश्यतावाचक मोडल तर्कशास्त्रात
$\varphi \strictif \psi$ ची व्याख्या
$\Box(\varphi \lif \psi)$ अशी करता येते.
म्हणून एखाद्या जगात कठोर अभिव्यंजन
सत्य असते तेव्हा आणि केवळ तेव्हाच संबंधित
वास्तविक अभिव्यंजन अवश्य सत्य असते.

वास्तविक अभिव्यंजनाच्या विरोधाभासांच्या
बाबतीत कठोर अभिव्यंजन कसे वागते?
पूर्वांग असत्य असलेले कठोर अभिव्यंजन,
तसेच उत्तरांग सत्य असलेले कठोर अभिव्यंजन,
सत्यही असू शकते किंवा असत्यही.
शिवाय $(\varphi \strictif \psi) \lor (\psi \strictif \varphi)$
वैध नाही. $\varphi \strictif \psi$
हे कठोर अभिव्यंजन
$\lnot \varphi \lor \psi$ शीही
सममूल्य नाही; त्यामुळे ते
सत्यता-फलनात्मक नाही.

S5 च्या संदर्भात पुढील संबंध मिळतात:
\begin{align}
  \varphi \strictif \psi & \Entails \lnot \varphi \lor \psi
  \text{ पण:}\\
  \lnot \varphi \lor \psi & \Entails/ \varphi \strictif \psi\\
  \psi & \Entails/ \varphi \strictif \psi\\
  \lnot \varphi & \Entails/ \varphi \strictif \psi\\
  \lnot(\varphi \strictif \psi) & \Entails/ \varphi \land \lnot \psi
  \text{ पण:}\\
  \varphi \land \lnot \psi & \Entails \lnot(\varphi \strictif \psi)
  \intertext{तरीही कठोर अभिव्यंजनामुळे
    विध्यात्मक अनुमान लागू होते:}
  \varphi, \varphi \strictif \psi & \Entails \psi
  \intertext{कठोर अभिव्यंजन
    निष्कर्षांचे संयुग्मीकरण करते:}
  \varphi \strictif \psi,  \varphi \strictif \chi & \Entails \varphi \strictif (\psi \land \chi)
  \intertext{कठोर अभिव्यंजनासाठी
    पूर्वांग अधिक बलवान करता येते:}
  \varphi \strictif \psi & \Entails (\varphi \land \chi) \strictif \psi
  \intertext{कठोर अभिव्यंजन
    संक्रमकही आहे:}
  \varphi \strictif \psi, \psi \strictif \chi & \Entails \varphi \strictif \chi
  \intertext{शेवटी, कठोर अभिव्यंजन
    त्याच्या प्रतिपरिवर्तित विधानाशी
    सममूल्य आहे:}
  \varphi \strictif \psi & \Entails \lnot \psi \strictif \lnot \varphi\\
  \lnot \psi \strictif \lnot \varphi & \Entails \varphi \strictif \psi
\end{align}

\begin{prob}
  कठोर अभिव्यंजनासाठी लागू न होणाऱ्या
  तार्किक निष्पन्नता-संबंधांना
  \Log{S5} मधील प्रतिदृष्टांत द्या;
  म्हणजे पुढील प्रत्येकासाठी:
  \begin{enumerate}
  \item $\lnot p \Entails/ \Box(p \lif q)$
  \item $q \Entails/ \Box(p \lif q)$
  \item $\lnot\Box(p \lif q) \Entails/ p \land \lnot q$
  \item $\Entails/ \Box(p \lif q) \lor \Box(q \lif p)$
  \end{enumerate}
\end{prob}

\begin{prob}
  वैध तार्किक निष्पन्नता-संबंध
  कठोर अभिव्यंजनासाठी लागू होतात
  हे दाखवण्यासाठी पुढील
  विधानांचे \Log{S5} मध्ये
  पुरावे द्या:
  \begin{enumerate}
  \item  $\Box(\varphi \lif \psi) \Entails \lnot \varphi \lor \psi$
  \item $\varphi \land \lnot \psi \Entails \lnot\Box(\varphi \lif \psi)$
  \item $\varphi, \Box(\varphi \lif \psi) \Entails \psi$
  \item $\Box(\varphi \lif \psi), \Box(\varphi \lif \chi) \Entails \Box(\varphi \lif (\psi
    \land \chi))$
  \item $\Box(\varphi \lif \psi) \Entails \Box((\varphi \land \chi) \lif \psi)$
  \item $\Box(\varphi \lif \psi), \Box(\psi \lif \chi) \Entails \Box(\varphi
    \lif \chi)$
  \item $\Box(\varphi \lif \psi) \Entails \Box(\lnot \psi \lif \lnot \varphi)$
  \item $\Box(\lnot \psi \lif \lnot \varphi) \Entails \Box(\varphi \lif \psi)$
  \end{enumerate}
  \end{prob}

तरी कठोर अभिव्यंजनाचे स्वतःचेही
‘‘विरोधाभास’’ आहेत. पूर्वांग असत्य
असेल किंवा उत्तरांग सत्य असेल,
तर वास्तविक अभिव्यंजन सत्य
असते. त्याचप्रमाणे, पूर्वांग
\emph{अवश्य} असत्य असेल
किंवा उत्तरांग अवश्य सत्य असेल,
तर कठोर अभिव्यंजन सत्य असते.
S5 मध्ये कोणतेही सत्य कठोर
अभिव्यंजन अवश्य सत्य असते
आणि कोणतेही असत्य कठोर
अभिव्यंजन अवश्य असत्य असते.
दुसऱ्या शब्दांत,
\begin{align}
  \Box \lnot \varphi & \Entails \varphi \strictif \psi\\
  \Box \psi & \Entails \varphi \strictif \psi\\
  \varphi \strictif \psi& \Entails \Box(\varphi \strictif \psi)\\
  \lnot(\varphi \strictif \psi) & \Entails \Box\lnot(\varphi \strictif \psi)
\end{align}
$\strictif$ चा अर्थ ‘‘तार्किकरीत्या
निष्पन्न होते’’ असा घेतल्यास
यात अडचण नाही. शेवटी,
तार्किक निष्पन्नतेचे संबंध ही
गणिती तथ्ये आहेत; म्हणून
ती नैमित्तिक असू शकत नाहीत.
पण ‘‘जर \dots तर \dots’’
याचा अर्थ पकडणारा तार्किक
कारक म्हणून $\strictif$
वापरायचा असेल, तर या
गोष्टी प्रश्न निर्माण करतात,
विशेषतः शेवटच्या दोन.
कारण काही ‘‘जर \dots तर \dots’’
विधाने नैमित्तिकरीत्या
सत्य किंवा असत्य असतात;
प्रत्यक्षात ती सहसा
ना अवश्य सत्य असतात,
ना अशक्य.

\begin{prob}
  पुढील विधानांचे \Log{S5}
  मध्ये पुरावे द्या:
  \begin{enumerate}
    \item  $\Box \lnot \varphi \Entails \varphi \strictif \psi$
    \item $\varphi \strictif \psi \Entails \Box(\varphi \strictif \psi)$
    \item $\lnot(\varphi \strictif \psi)  \Entails \Box\lnot(\varphi \strictif \psi)$
  \end{enumerate}
  त्यासाठी $\strictif$ ची
  व्याख्या वापरा.
\end{prob}












\section{प्रतिवास्तविक विधाने}\label{cnt:int:cnt:sec}

इंग्रजीतील ‘‘जर \dots तर \dots’’
या रचनेचा एक अतिशय प्रचलित आणि
महत्त्वाचा प्रकार \emph{to be}
या क्रियापदाच्या भूतकालीन
संकेतार्थी रूपाने बनतो:
‘‘जर असे घडले असते की \dots,
तर असे घडले असते की \dots’’.
अशा सशर्त विधानाचे पूर्वांग
सहसा असत्य, म्हणजे वस्तुस्थितीच्या
विरुद्ध असल्याने, त्यांना
\emph{प्रतिवास्तविक सशर्त विधाने}
म्हणतात. \emph{to be} चे
संकेतार्थी रूप त्यांत असल्याने
त्यांना \emph{संकेतार्थी सशर्त विधाने}
असेही म्हणतात. याउलट
\emph{निश्चयार्थी} सशर्त विधाने
‘‘जर असे आहे की \dots,
तर असे आहे की \dots’’
या रूपात असतात. प्रतिवास्तविक
आणि निश्चयार्थी सशर्त विधानांच्या
सत्यतेच्या अटी वेगळ्या असतात.
अॅडम्स यांचे प्रसिद्ध उदाहरण पाहा:
\begin{quote}
  जर ऑस्वाल्डने केनेडींची हत्या केली नाही,
  तर दुसऱ्या कोणीतरी केली.

  जर ऑस्वाल्डने केनेडींची हत्या केली
  नसती, तर दुसऱ्या कोणीतरी केली असती.
\end{quote}
पहिले विधान निश्चयार्थी, तर दुसरे
प्रतिवास्तविक आहे. पहिले विधान
स्पष्टपणे सत्य आहे: अध्यक्ष जॉन
एफ. केनेडी यांची हत्या
\emph{कोणीतरी} केली हे आपल्याला
माहीत आहे. ती व्यक्ती,
वॉरेन अहवालात म्हटल्याच्या विरुद्ध,
ली हार्वे ऑस्वाल्ड नसेल,
तर दुसऱ्या कोणीतरी हत्या
केली. दुसरे विधान मात्र
वेगळे काही सांगते. त्याचा
दावा असा की ऑस्वाल्डने
केनेडींची हत्या केली नसती
म्हणजे डॅलस येथील गोळीबार
टळला असता किंवा अयशस्वी
ठरला असता, तरी पुढचा
इतिहास अशा प्रकारे घडला
असता की दुसरा हत्येचा
प्रयत्न यशस्वी झाला असता.
हे प्रतिवास्तविक विधान सत्य असण्यासाठी, ऑस्वाल्डने
हत्या केली नसती तरी दुसऱ्याने हत्या केली असती,
अशी त्या पर्यायी परिस्थितीविषयीची अट आवश्यक आहे.
मृत्यू निश्चित व्हावा म्हणून काही प्रबळ शक्तींनी
कट रचलेला असणे हे तिचे एक संभाव्य उदाहरण आहे
(अनेक कटसिद्धांतवादी जसे मानतात तसे); ते एकमेव
संभाव्य कारण असल्याचा तार्किक निष्कर्ष मात्र निघत नाही.

\emph{निश्चयार्थी} सशर्त विधान
वास्तविक अभिव्यंजनाने अचूक
दर्शवले जाते का, हा अजून
वादाचा विषय आहे. विशेषतः,
इंग्रजी निश्चयार्थी सशर्त
विधानांच्या सत्यतेच्या अटी
वास्तविक अभिव्यंजन देते
या मताशी सुसंगत राहील
अशा प्रकारे त्याचे
विरोधाभास ‘‘समजावता’’
येतात का, हा प्रश्न
आहे. याउलट प्रतिवास्तविक
सशर्त विधाने वास्तविक
अभिव्यंजनाने अचूक
दर्शवता येत नाहीत,
यावर वाद नाही.
कारण प्रतिवास्तविक
विधानांची पूर्वांगे
सहसा असत्य असली,
तरी असत्य पूर्वांग
असलेले प्रत्येक
प्रतिवास्तविक विधान
सत्य नसते. उदाहरणार्थ,
ऑस्वाल्डने हत्या केली नसती तर इतर कोणाकडूनही
हत्या झाली नसती, अशी पर्यायी परिस्थिती गृहीत धरल्यास,
अॅडम्स यांचे प्रतिवास्तविक विधान असत्य ठरते.
वॉरेन अहवाल खरा मानणे आणि कट नसणे या दोन
गृहीतकांवरूनच ही अतिरिक्त प्रतिवास्तविक अट
तार्किकरीत्या निष्पन्न होते, असा दावा येथे करत नाही.

प्रतिवास्तविक सशर्त विधाने
कार्यकारणाविषयीच्या
युक्तिवादात महत्त्वाची
असतात: प्रतिवास्तविक
विधानांनी कारणसंबंध
दर्शवणे हे त्याचे
प्रमुख उदाहरण आहे.
उदाहरणार्थ, आगकाडी
ओढल्याने ती पेटते;
हे ‘‘ही आगकाडी
ओढली असती, तर ती
पेटली असती’’ असे
म्हणून दर्शवता येते.
वास्तविक अभिव्यंजनाने,
किंवा सर्वसाधारणपणे
निश्चयार्थी सशर्त विधानांनी,
हे दर्शवता येत नाही:
‘‘आगकाडी ओढली जाते
$\lif$ आगकाडी पेटते’’
हे विधान आगकाडी
कधीच ओढली गेली
नाही तर सत्य असते,
ती ओढली असती
तर काय झाले असते
याचा काहीही संबंध
नसतो. त्याहूनही
वाईट म्हणजे,
‘‘आगकाडी ओढली जाते
$\lif$ आगकाडीचे
फुलांच्या गुच्छात
रूपांतर होते’’ हेही
आगकाडी कधीच
ओढली गेली नाही
तर सत्य असते;
पण ती ओढली
असती तर नक्कीच
तिचे फुलांच्या
गुच्छात रूपांतर
झाले नसते.

\begin{quote}\small\textbf{स्रोतनोंद.} OLINC-238; SOL6-A212: एका सर्वाधिक जवळच्या जगाची मूळ सामायिक सत्यता-अट दुरुस्त केली आहे. स्टालनाकरच्या निवडलेल्या जगाची अट आणि लुईसची सर्व पूर्वांग-जगांना लागू होणारी गोलक-अट वेगळी दिली आहे; सर्वात जवळचे जग अस्तित्वातच नसेल तरी लुईसची अट लागू होते. मागील हत्या-उदाहरणातील कट हा एकमेव आवश्यक कारण असल्याचा अतिरिक्त दावा काढला आहे.\end{quote}
प्रतिवास्तविक विधानांचे योग्य तर्कशास्त्र नेमके काय,
यावर अजूनही वाद आहे. स्टालनाकर आणि लुईस यांनी
प्रभावी, पण वेगवेगळी विश्लेषणे दिली.
‘‘जर~$S$ असे घडले असते, तर~$T$ असे घडले असते’’
या विधानासाठी, स्टालनाकरच्या मांडणीत~$S$ सत्य
असलेले प्राप्य जग उपलब्ध असल्यास एक जवळचे जग
निवडले जाते; त्या निवडलेल्या जगात~$T$ सत्य असेल
तेव्हा आणि केवळ तेव्हाच प्रतिवास्तविक विधान सत्य असते.
अशक्य पूर्वांगासाठी त्याच्या औपचारिक मांडणीत
विशेष असंभव जग वापरले जाते.
लुईसची अट एका अद्वितीय जवळच्या जगावर अवलंबून नाही:
एकतर संबंधित गोलकांत~$S$ सत्य असलेले जगच नसते,
किंवा~$S$ सत्य असलेले किमान एक जग समाविष्ट करणारा
असा गोलक असतो की त्यातील प्रत्येक $S$-जगात~$T$
सत्य असते. सर्वाधिक जवळची $S$-जगे अस्तित्वात
असतील, तर ही अट त्या सर्व जगांत~$T$ सत्य असण्याशी
सममूल्य आहे; त्यांत अनेक जगे असू शकतात.
सर्वाधिक जवळचे $S$-जग नसतानाही गोलक-अट लागू होते.
हे भेद स्टालनाकर यांच्या मूळ लेखाच्या विभाग~II मध्ये
आणि लुईस यांच्या मूळ ग्रंथातील विभाग~1.3 मध्ये दिले आहेत
(\url{https://philosophy.hku.hk/courses/dm/phil2520/ATheoryOfConditionals-Stalnaker.pdf};
\url{https://fitelson.org/piksi/piksi_18/lewis_counterfactuals.pdf}, पृ.~16, 20).
संभाव्य जगांच्या
सत्ताशास्त्राचा
संदर्भ घेत असल्याने
याला ‘‘सत्ताविषयक’’
विश्लेषण म्हणतात.
इतर विश्लेषणे सशर्त
संभाव्यता किंवा
विश्वास-पुनरीक्षणाचे
सिद्धांत वापरतात.
प्रतिवास्तविक विधानांसाठी
अनेक वेगवेगळी
तर्कशास्त्रे सुचवली
गेली आहेत. लुईस--स्टालनाकर
यांचे एकच तर्कशास्त्रही
नाही: सर्वाधिक जवळच्या
संभाव्य जगांच्या आधारे
त्यांनी साधर्म्य असलेली
मांडणी केली, तरी
त्यांच्या भिन्न
गृहीतकांमुळे भिन्न
वैध अनुमाने मिळतात.



